%
% --- color support (needed by \textcolor in \PH and \definecolor in main.tex)
%
\usepackage{xcolor}

%
% --- inline annotations
%
\newcommand{\red}[1]{{\color{red}#1}}
\newcommand{\todo}[1]{{\color{red}#1}}
\newcommand{\TODO}[1]{\textbf{\color{red}[TODO: #1]}}
% --- disable by uncommenting
% \renewcommand{\TODO}[1]{}
% \renewcommand{\todo}[1]{#1}

% --- placeholder for anything not in the evidence packet
\newcommand{\PH}[1]{\textcolor{red}{[PH:#1]}}

% --- figures live in figures/
\usepackage{graphicx}
\usepackage{subcaption}
\graphicspath{{figures/}}

% --- math macros
\usepackage{amsmath,amssymb,amsthm}
\setlength{\emergencystretch}{3em} % absorb stubborn overfull lines during page build
% Two-column article defaults to \flushbottom, which stretches sparse columns and
% opens large inter-paragraph gaps (especially around the nobreak result/proposition
% boxes). Let column bottoms be ragged instead of vertically stretched.
\raggedbottom

% The official review submission displays WACV's line numbers.

% --- framed "Result" boxes (Allen-Zhu "Physics of LM" style); pdflatex-safe
\usepackage{mdframed}
\usepackage{booktabs}  % \toprule/\midrule/\bottomrule in body curated tables
% Proxy/metric names: clean sans-serif identifier (NOT a monospace "code" look),
% hyphenated to match the figure labels, and breakable to avoid column overflow.
\newcommand{\px}[1]{\textsf{#1}}
\newcommand{\E}{\mathbb{E}}
\newcommand{\Var}{\operatorname{Var}}
\newcommand{\Corr}{\operatorname{Corr}}
\newcommand{\Phimc}{\mathcal{G}}
\DeclareMathOperator{\erank}{erank}
\newcommand{\corr}{\operatorname{corr}}
\newcommand{\partialcorr}{\operatorname{partial}}
\newcommand{\acc}{\mathrm{acc}}
\newcommand{\params}{\mathrm{params}}
\theoremstyle{plain}
\newtheorem{theorem}{Theorem}
\newtheorem{proposition}{Proposition}
\theoremstyle{definition}
\newtheorem{definition}{Definition}
% --- empirical, falsifiable regularities (NOT theorems): numbered "Result N"
% Rendered as a light-gray framed box (Allen-Zhu "Result N" register).
\theoremstyle{definition}
\newtheorem{findingthm}{Result}
\mdfdefinestyle{findingbox}{%
  backgroundcolor=black!4,
  linecolor=black!30,
  linewidth=0.4pt,
  roundcorner=2pt,
  innertopmargin=3pt,
  innerbottommargin=3pt,
  innerleftmargin=4pt,
  innerrightmargin=4pt,
  skipabove=3pt,
  skipbelow=3pt,
  nobreak=false, % allow the box to split across a column/page break, so a box that
                 % does not fit at a column bottom FILLS it and continues, instead of
                 % being bumped whole and leaving a large ragged gap below.
}
\newmdenv[style=findingbox]{findingframe}
\newenvironment{finding}[1][]%
  {\begin{findingframe}\begin{findingthm}[#1]}%
  {\end{findingthm}\end{findingframe}}
% --- inline run-in numbered "Result" environment, sharing the findingthm counter.
% Renders "\textbf{Result k (handle).} ..." inline; \cref-able as "Result k".
\newtheoremstyle{resultrunin}{4pt}{4pt}{}{}{\bfseries}{.}{ }{\thmname{#1}~\thmnumber{#2}\thmnote{ (#3)}}
\theoremstyle{resultrunin}
\newtheorem{result}[findingthm]{Result}
% cleveref names for \cref{find:...} (cleveref is loaded \AtEndPreamble in wacv.sty)
% The counter driving \cref is now `findingthm` (the inner theorem inside the box).
\AtEndPreamble{%
  \crefname{findingthm}{Result}{Results}%
  \Crefname{findingthm}{Result}{Results}%
}

